\documentclass[aspectration=1610,t]{beamer}
\usepackage{csc}
\title{Лекция 1. Вводная}


\date{
   \textbf{CS центр}\\
   7 сентября 2021\\
   Санкт-Петербург
}

\begin{document}

\begin{frame}
  \titlepage
\end{frame}

\begin{frame}[fragile]{Некоторые вводные}
    Ожидается, что люди, пришедшие на курс:
    \begin{itemize}
        \item Имеют опыт программирования на любом языке программирования.
        \item Имеют базовые представления об ООП:
            классы, объекты, икапсуляция, наследование и полиморфизм.
        \item Знакомы с базовыми алгоритмами и структурами данных: 
            поиск, сортировка, линейный массив, хеш-таблица, \dots
    \end{itemize}
\end{frame}

\begin{frame}[fragile]{Какие результаты ожидаем}
    \begin{itemize}
        \item Сможем самостоятельно решать задачи с использованием \langcpp.
        \item Научимся самостоятельно находить ответы на свои вопросы.
        \item Будем писать эффективный, простой и быстрый код,
            применяя естественные для \langcpp подходы.
    \end{itemize}
\end{frame}

\begin{frame}[fragile]{Где искать ответы на вопросы}
    Опишем примерный алгоритм поиска ответы на вопросы:
    \begin{itemize}
        \item Гуглим проблему и идем на \url{https://stackoverflow.com}.
        \item Читаем человекочитаемую документацию на 
            \url{https://en.cppreference.com/w/}.
        \item Смотрим видеозаписи курсов на \url{https://stepik.org/course/7/} и 
            \url{https://stepik.org/course/3206/}.
        \item Читаем актуальную версию черновика стандарта 
            \url{https://github.com/cplusplus/draft}.
    \end{itemize}
\end{frame}

\begin{frame}{Особенности \langc}
    \begin{itemize}
        \item {\bf Эффективность.}\\
            Язык С позволяет писать программы, которые напрямую работают с железом. 
        \item {\bf Стандартизированность.}\\
            Спецификация языка \langc является международным стандартом.
        \item {\bf Относительная простота.}\\
            Стандарт языка \langc занимает 230 страниц\\
            (против 700+ для \texttt{Java} и 1300+ для \langcpp).
    \end{itemize}
\end{frame}

\begin{frame}{Совместимость \langc и \langcpp}

    \begin{itemize}
        \item Один из принципов разработки стандарта \langcpp~— это
            cохранение совместимости с \langc. 
        \item Синтаксис \langcpp унаследован от языка \langc. 
        \item \langcpp не является в строгом смысле надмножеством \langc.
        \item Можно писать программы на \langc так, чтобы они успешно
            компилировались на \langcpp.
        \item \langc и \langcpp сильно отличаются как по сложности, так и по принятым 
            архитектурным решениям, которые используются в обоих языках.
    \end{itemize}
\end{frame}

\begin{frame}{Стандартизация \langcpp}
    \begin{itemize}
        \item Лишь в 1998 году был ратифицирован международный стандарт языка 
            \langcpp: ISO/IEC 14882:1998 ``Standard for the \langcpp Programming
            Language''.
        \item В 2003 году был опубликован стандарт языка ISO/IEC 14882:2003, где были
            исправлены выявленные ошибки и недочёты предыдущей версии стандарта.
        \item С 2005 года началась работа над новой версией стандарта, которая
            получила кодовое название C++0x. В конце концов в 2011 году стандарт 
            был принят и получил название \langcpp[11] ISO/IEC 14882:2011. 
        \item В 2014 году вышел \langcpp[14]: ISO/IEC 14882:2014.
        \item В 2017 году вышел \langcpp[17]: ISO/IEC 14882:2017.
        \item В 2020 году вышел \langcpp[20]: ISO/IEC 14882:2020.
        \item В данный момент готовится к публикации \langcpp[23].
    \end{itemize}
\end{frame}

\begin{frame}{Характеристики языка \langcpp}
    Характеристики \langcpp:
    \begin{itemize}
        \item сложный,
        \item мультипарадигмальный,
        \item эффективный,
        \item низкоуровневый,
        \item компилируемый,
        \item статически типизированный.
    \end{itemize}
\end{frame}

\begin{frame}[fragile]{Сложность}
    \begin{itemize}
        \item Описание стандарта занимает более 1300 страниц текста.
        \item Нет никакой возможности рассказать ``весь \langcpp'' в рамках
            одного, пусть даже очень большого курса.
        \item В \langcpp программисту позволено очень многое, и это влeчёт за
            собой большую ответственность.
        \item На плечи программиста ложится много дополнительной работы:
            \begin{itemize}
                \item проверка корректности данных, 
                \item управление памятью, 
                \item обработка низкоуровневых ошибок,
                \item \dots
            \end{itemize}
    \end{itemize}
\end{frame}

\begin{frame}{Мультипарадигмальный}
    На \langcpp можно писать программы в рамках нескольких парадигм программирования:
    \begin{itemize}
        \item {\bf процедурное программирование}\\
            (код ``в стиле \langc''),
        \item {\bf объектно-ориентированное программирование}\\
            (классы, наследование, виртуальные функции, \dots).
        \item {\bf обобщённое программирование}\\
            (шаблоны функций и классов),
        \item {\bf функциональное программирование}\\
            (функторы, безымянные функции, замыкания),
        \item {\bf генеративное программирование}\\
            (метапрограммирование на шаблонах).
    \end{itemize}
\end{frame}

\begin{frame}{Эффективный}
    Одна из фундаментальных идей языков \langc и \langcpp~---
    {\em отсутствие неявных накладных расходов}, которые присутствуют
    в других более высокоуровневых языках программирования.
    \begin{itemize}
        \item Программист сам выбирает уровень абстракции, на котором писать
            каждую отдельную часть программы. 
        \item Можно реализовывать критические по производительности участки
        программы максимально эффективно.
        \item Эффективность делает \langcpp основным языком для разработки 
            высоконагруженных приложений, приложений с компьютерной графикой и \dots
    \end{itemize}

\end{frame}

\begin{frame}{Низкоуровневый}
    Язык \langcpp, как и \langc, позволяет работать напрямую 
    с ресурсами компьютера.
    \begin{itemize}
        \item Позволяет писать низкоуровневые системные приложения 
            (например, драйверы операционной системы).
        \item Неаккуратное обращение с системными ресурсами 
            может привести к падению программы.
    \end{itemize}

    В \langcpp отсутствует автоматическое управление памятью. 
    \begin{itemize}
        \item Позволяет программисту получить полный контроль над программой.
        \item Необходимость заботиться об освобождении памяти.
    \end{itemize}
\end{frame}

\begin{frame}{Компилируемый}
    \langcpp является компилируемым языком программирования.
    \medskip

    Для того, чтобы запустить программу на \langcpp, её нужно сначала {\em скомпилировать}.
    \medskip

    Компиляция~--- преобразование текста программы на языке программирования
    в машинный код.

    \begin{itemize}
        \item Нет накладных расходов при исполнении программы.
        \item При компиляции можно отловить некоторые ошибки.
        \item Требуется компилировать для каждой платформы отдельно.
    \end{itemize}

\end{frame}

\begin{frame}{Статическая типизация}
    \langcpp является статически типизированным языком.
    \medskip

    \begin{enumerate}
        \item Каждая сущность в программе (переменная, функция и пр.) имеет свой тип,
        \item и этот тип определяется на момент компиляции. 
    \end{enumerate}
    Это нужно для того, чтобы
    \begin{enumerate}
        \item вычислить размер памяти, который будет занимать каждая переменная в программе,
        \item определить, какая функция будет вызываться в каждом конкретном месте. 
    \end{enumerate}

    Всё это определяется на момент компиляции и ``зашивается'' в скомпилированную программу. 

    В машинном коде никаких типов уже нет~--- там идёт работа с последовательностями байт. 
\end{frame}


\begin{frame}[fragile]{Что такое компиляция?}
\begin{center}
    \begin{tikzpicture}[auto,thick,font=\tiny]
    \tikzstyle{arrow} = [>=stealth',shorten >=2pt, shorten <=2pt]
    \tikzstyle{file} = [rectangle,minimum width=7ex,minimum height=15mm,draw,thick]

    \node [shift={(0mm,0mm)}] (programmer) {\includegraphics[height=1cm]{programmer}};
    \node [file,shift={(35mm,0cm)}]  (file1) {\parbox{14mm}{Архитектура приложения}};
    \node [file,shift={(80mm,0cm)}]  (file2) {Код на \langcpp};
    \node [file,shift={(35mm,-3cm)}] (file3) {Машинный код};
    \node [shift={(80mm,-3cm)}] (computer) {\includegraphics[height=1cm]{computer}};
    \path [arrow,->] (programmer) edge node {Проектирование} (file1);
    \path [arrow,->] (file1) edge node {Программирование}  (file2);
    \path [arrow,->] (file3) edge node {Исполнение}        (computer);

    \draw[->, arrow] (file2.south) -- ++(0,-.6) -| (file3.north);
    \node [shift={(57mm,-12mm)}] {Компиляция};
    
    \node[cloud, cloud puffs=15.7, cloud ignores aspect, align=center, draw]
        (cloud1) at (0cm, 1cm) {Идея};
    \node[cloud, cloud puffs=15.7, cloud ignores aspect, align=center, draw]
        (cloud2) at (80mm, -2cm) {0101};
\end{tikzpicture}
\end{center}
\end{frame}

\begin{frame}[fragile]{Что такое компиляция?}
\begin{center}
\begin{tikzpicture}[auto,thick,font=\tiny]
    \tikzstyle{arrow} = [>=stealth',shorten >=2pt, shorten <=2pt]
    \tikzstyle{file} = [rectangle,minimum width=7ex,minimum height=15mm,draw,thick]

    \node [shift={(0mm,0mm)}] (programmer) {\includegraphics[height=1cm]{programmer}};
    \node [file,shift={(35mm,0cm)}]  (file1) {\parbox{14mm}{Архитектура приложения}};
    \node [file,shift={(80mm,0cm)}]  (file2) {Код на \texttt{Java}};
    \node [file,shift={(0mm,-3cm)}] (file3) {Байт код};
    \node [cloud, draw, shift={(35mm,-3cm)},label=below:JVM] (vm) {\includegraphics[height=8mm]{computer}};
    \node [shift={(80mm,-3cm)}] (computer) {\includegraphics[height=1cm]{computer}};
    \path [arrow,->] (programmer) edge node {Проектирование} (file1);
    \path [arrow,->] (file1) edge node {Программирование}  (file2);


    \draw [arrow,->] (file2.south) -- ++(0,-.6) -| (file3.north);
    \node [shift={(35mm,-12mm)}] {Компиляция};

    \path [arrow,->] (file3) edge node {Исполнение}        (vm);
    \path [arrow,<->] (vm) edge node {Трансляция команд}          (computer);
    
    \node[cloud, cloud puffs=15.7, cloud ignores aspect, align=center, draw]
        (cloud1) at (0cm, 1cm) {Идея};
    \node[cloud, cloud puffs=15.7, cloud ignores aspect, align=center, draw]
        (cloud2) at (80mm, -2cm) {0101};
\end{tikzpicture}
\end{center}
\end{frame}

\newcommand{\gear}[5]{%
\foreach \i in {1,...,#1} {%
  [rotate=(\i-1)*360/#1]  (0:#2)  arc (0:#4:#2) {%[rounded corners=1.5pt]
               -- (#4+#5:#3)  arc (#4+#5:360/#1-#5:#3)} --  (360/#1:#2)
               }} 

\begin{frame}[fragile]{Что такое интерпретация?}
\begin{center}
\begin{tikzpicture}[auto,thick,font=\tiny]
    \tikzstyle{arrow} = [auto, thick, >=stealth',shorten >=2pt, shorten <=2pt]
    \tikzstyle{file} = [rectangle,minimum width=1cm,minimum height=15mm,draw,thick]

    \node [shift={(0mm,0mm)}] (programmer) {\includegraphics[height=1cm]{programmer}};
    \node [file,shift={(35mm,0cm)}]  (file1) {\parbox{14mm}{Архитектура приложения}};
    \node [file,shift={(80mm,0cm)}]  (file2) {Код на \texttt{Perl}};
    \node [shift={(35mm,-3cm)},label=below:Интерпретатор,minimum
    height=15mm,minimum width=15mm] (perl) {Perl};
    
    \node [shift={(80mm,-3cm)}] (computer) {\includegraphics[height=1cm]{computer}};
    \path [arrow,->] (programmer) edge node {Проектирование} (file1);
    \path [arrow,->] (file1) edge node {Программирование}  (file2);
    \path [arrow,<->] (perl)   edge node {Трансляция команд} (computer);

    \draw [->, arrow] (file2.south) -- ++(0,-.6) -| (perl.north);
    \node [shift={(57mm,-12mm)}] {Интерпретация};
    
    \draw [thick,shift={(35mm,-3cm)}] \gear{10}{0.6}{.7}{10}{2};
    \draw [thick,shift={(35mm,-3cm)}] circle (.4); 
                                     
    \node [cloud, cloud puffs=15.7, cloud ignores aspect, align=center, draw]
        (cloud1) at (0cm, 1cm) {Идея};
    \node [cloud, cloud puffs=15.7, cloud ignores aspect, align=center, draw]
        (cloud2) at (80mm, -2cm) {0101};
\end{tikzpicture}
\end{center}
\end{frame}

%С одной стороны находится человек, с другой -- компьютер. Компьютер и человек
%говорят на совершенно разных языках. Программирование -- это процесс перевода
%желаний человека на язык, понятный компьютеру.
%
%Это преобразование занимает несколько этапов.
%Предположим, что человек хочет написать графический редактор.
%Он понимает, как он работать с картинками.
%
%Для упрощения процесса будем считать, что человек -- это программист.
%Взаимодействие заказчика и программиста -- это отдельный сложный вопрос, который
%выходит за границы данного курса.
%
%Первый этап -- проектирование и алгоритмизация. На этом этапе программист
%описывает устройство программы в терминах сущностей, которые предоставляет
%конкретный язык программирования. К примеру, если мы говорит про ООП, то
%программист выделяет подсистемы, каждую из которых разбивает на классы, 
%выделяет интерфейсы классов и описывает объектную модель взаимодействия 
%между ними. Для каждой операции выбирается конкретный алгоритм, к примеру,
%для преобразования цветного изображения в черно-белое можно выбрать алгоритм,
%которые преобразует яркость пикселя в соответствующее значение ч/б пикселя 
%компоненту.
%
%Второй этап -- это реализация предложенного архитектурного и алгоритмического
%решения в терминах конкретного языка программирования. Т.е. в нашем примере, это
%непосредственное кодирование предложенной объектной модели приложения и
%выбранных алгоритмов на каком-то языке программирования. Результатом данного
%этапа является код программы выбранном на языке программирования.
%
%Третий этап -- это компиляция. В этот момент программа из языка, понятном
%программисту, транслируется на какой-то другой язык. Обычно в данном случае
%подразумевается язык понятный компьютеру -- в машинные коды.
%Иногда компиляция может происходить в несколько этапов. Так, к примеру, первый
%компилятор С++ cfront компилировал код на С++ в код на С, а потом уже вызывался
%компилятор С, который переводил код в машинные коды. Так и сейчас работают
%некоторые компиляторы С++ -- к примеру, компилятор Comeau. Язык С в данном
%случая, как я уже говорил, выступает в роли универсального ассемблера -- языка
%низкого уровня, который просто транслировать в машинные коды конкретной
%архитектуры. Существует множество языков, которые компилируются в С: Scheme,
%Haskell, OCaml и т.д.
%
%Очень часто язык программирования компилируется в язык ассемблера -- это такой
%специальный язык, который очень близок по структуре к машинным кодам конкретной
%архитектуры. Компилирование программы на языке ассемблера (ассемблирование) --
%это очень простой процесс, который можно представить как замену по словарю,
%когда каждой команде (к примеру, ``mov (%eax),%eax'' превратится в два байта
%``8b00'').
%
%Есть ещё и другие подходы. Антонимом к ``компилируемому языку'' является
%``интерпретируемый язык''. В этом случае этап 3 отсутствует, а вместо него
%используется специальная программа -- интерпретатор, которая читает код на языке
%программирования и выполняет все операции строчка за строчкой. По такой схеме
%работают различные скриптовые языки, к примеру, PHP и JavaScript.
%
%Ещё один вариант -- использование виртуальной машины. Для таких языков на этапе
%3 генерируется код не для конкретного процессора, а для некоторой абстрактного
%не существующего в природе процессора, который на каждой архитектуре реализован 
%в виде специальной программы -- виртуальной машины. Данный подход -- нечто 
%среднее между компилируемым языком и интерпретируемым, позволяющий сочетать 
%некоторые плюсы обоих решений. Так работают языки семейств Java и .Net.
%
%Давайте закончим данный фрагмент перечислением плюсов и минусов
%компилируемости языка.

\begin{frame}[fragile]{Плюсы и минусы компилируемости в машинный код}{}
    {\bf Плюсы:}
    \begin{itemize}
        \item эффективность: программа компилируется и оптимизируется для
            конкретного процессора,
        \item нет необходимости устанавливать сторонние приложения (такие как
            интерпретатор или виртуальная машина).
    \end{itemize}

    {\bf Минусы:}
    \begin{itemize}
        \item нужно компилировать для каждой платформы,
        \item сложность внесения изменения в программу~--- 
            нужно перекомпилировать заново.
    \end{itemize}

    {\bf Важно:} компиляция — преобразование одностороннее, нельзя восстановить исходный код.
\end{frame}

\begin{frame}[fragile]{Разбиение программы}
    Зачем разбивать программу на файлы? 
    \begin{itemize}
        \item С небольшими файлами удобнее работать.
        \item Разбиение на файлы структурирует код.
        \item Позволяет нескольким программистам разрабатывать 
            приложение одновременно.
        \item Ускоряет компиляцию программы.
        \item Ускорение повторной компиляции при небольших 
            изменениях в отдельных частях программы. 
            
    \end{itemize}
        Файлы с кодом на \langcpp бывают двух типов:
        \begin{enumerate}
            \item файлы с исходным кодом (расширение \texttt{.cpp}, иногда \texttt{.C}),
            \item заголовочные файлы (расширение \texttt{.hpp} или \texttt{.h}).
        \end{enumerate}
\end{frame}

%Для того, чтобы говорить о компиляции С++ кода нужно сначала изучить
%его структуру. Программа на С++ состоит из двух типов файлов:
%файлов с исходным кодом (обычно с расширением .cpp или .С) и заголовочных файлов
%(с расширением .h или .hpp). Предполагается, что весь код написан в файлах с
%исходным кодом, а заголовочные файлы служат только для описания зависимостей.
%
%Давайте ответим на простой вопрос: зачем разбивать программу на
%файлы? Ответ просто: потому, что нам удобнее работать с небольшими файлами,
%чем с одним огромным, в котором ничего не найти. Это так же позволяет
%разрабатывать приложение по частям параллельно, когда каждый программист
%работает над какими-то своими файлами. Есть и другие причины для разделения на
%файлы -- это переиспользование кода (в случае библиотек) и ускорение перекомпиляции
%при небольших изменениях в отдельных частях программы. В таких случаях, если
%изменения затронули небольшое количество файлов, достаточно перекомпилировать
%только их.

\begin{frame}[fragile]{Простейшая программа на С++}
    \begin{itemize}
        \item Файл {\tt foo.hpp}:
            {\small \begin{lstlisting}
#pragma once // страж включения

// объявление (declaration) функции foo
void foo();
            \end{lstlisting}}

        \item Файл {\tt foo.cpp}:
            \begin{lstlisting}
#include "foo.hpp"
#include "bar.hpp"

// определение (definition) функции foo
void foo()
{
    bar(); // вызываем функцию bar
}
            \end{lstlisting}
    \end{itemize}
\end{frame}

\begin{frame}[fragile]{Простейшая программа на С++}
    \begin{itemize}
        \item Файл {\tt bar.hpp}:
            {\small \begin{lstlisting}
#pragma once // страж включения

// объявление (declaration) функции bar
void bar();
            \end{lstlisting}}

        \item Файл {\tt bar.cpp}:
            \begin{lstlisting}
#include "bar.hpp"

// определение (definition) функции bar
void bar()
{
    // полезный код
}
            \end{lstlisting}
    \end{itemize}
\end{frame}

\begin{frame}[fragile]{Простейшая программа на С++}
    \begin{itemize}
        \item \emph{Точка входа} — функция, вызываемая при запуске программы.
            По умолчанию — это функция \verb!main!:
            {\small \begin{lstlisting}
#include "foo.hpp"

int main()
{
    foo();
}
            \end{lstlisting}}
            или
            {\small \begin{lstlisting}
#include "foo.hpp"

int main(int argc, char ** argv)
{
    foo();
}
            \end{lstlisting}}
    \end{itemize}
\end{frame}

\begin{frame}[fragile]{Замечания}
    \begin{itemize}
        \item Не забывайте разбивать программу на части.
        \item Не забывайте про стражи включения.
        \item Отличайте определения (.cpp) от объявлений (.hpp).
    \end{itemize}
\end{frame}

\begin{frame}[fragile] {Компилиция этап №1: препроцессор}
    \begin{itemize}
        \item Язык препроцессора -- это специальный язык программирования, 
            встроенный в \langcpp.
        \item Препроцессор работает с кодом на \langcpp как с текстом.
        \item Команды языка препроцессор называют директивами,              
            все директивы начинаются со знака \code{\#}.
        \item Директива \code{\#include} позволяет подключать заголовочные файлы к файлам
            кода.
        \begin{enumerate}
            \item \code{\#include <foo.h>} --- библиотечный заголовочный файл,
            \item \code{\#include "bar.h"{}} --- локальный заголовочный файл.
        \end{enumerate}
        \item Препроцессор заменяет директиву \code{\#include "bar.h"{}} на содержимое
        файла \texttt{bar.h}.
    \end{itemize}

%Тут важно отметить, что язык препроцессора ничего не знает о коде на С++. Т.е.
%потенциально можно использовать препроцессор С++ для других языков, т.к. после
%работы препроцессора всего его директивы из файлов удаляются. Это важно понимать
%потому, что иногда программисты стараются использовать конструкции препроцессора
%(макросы) для реализации некоторых конструкций языка (к примеру, функций). Это
%может привести к очень нехорошим последствиям. Мы это обсудим в дальнейшем.
%
%Чтобы посмотреть, что получилось после работы препроцессора нужно при компиляции
%указать ключ -E: 
%g++ -E program.cpp > program_preprocessed.cpp

\end{frame}

\begin{frame}[fragile]{Компилиция этап  2: компиляция}
    \begin{itemize}
        \item На вход компилятору поступает код на \langcpp после обработки
            препроцессором.
        \item Каждый файл с кодом компилируется отдельно и независимо от других
            файлов с кодом.
        \item Компилируется только файлы с кодом (т.е. {\tt *.cpp}).
        \item Заголовочные файлы сами по себе ни во что не компилируются, только в
            составе файлов с кодом.
        \item На выходе компилятора из каждого файла с кодом получается ``объектный
            файл''~--- бинарный файл со скомпилированным кодом (с расширением
            \texttt{.o} или \texttt{.obj}).
    \end{itemize}
\end{frame}
%Чтобы скомпилировать файл в объектный код нужно указать ключ -с 
%(g++ -с program.cpp).
%При желании можно попросить компилятор транслировать код в ассемблер
%(g++ -S program.cpp).

\begin{frame}[fragile]{Компилиция этап 3: линковка (компоновка)}
    \begin{itemize}
        \item На этом этапе все объектные файлы объединяются в один исполняемый 
            (или библиотечный) файл.
        \item При этом происходит подстановка адресов функций в места их вызова.
        \begin{lstlisting}
    void foo()
    {
        bar();
    }
    \end{lstlisting}
        \item По каждому объектному файлу строится таблица всех
            функций, которые в нём определены.
    \end{itemize}
\end{frame}
%В данном примере в теле функции foo вызывается функция bar. Этой функции нет в
%файле foo.cpp. По этому при компиляции в место вызова функции bar будет
%вставлено указание "вызвать функцию bar". При линковке компоновщик
%просмотрит все объектные файлы и составит список всех объявленных функций. Потом
%он заменит каждое указание вида "вызвать функцию ..." на инструкцию вызова
%функции: посмотрит в таблице всех функций и если найдёт такую, то заменит эту
%инструкцию на вызов функции по адресу. Если такой инструкции не будет --
%произойдёт ошибка. Если вдруг окажется, что есть несколько одинаковых функций в
%разных объектных файлах, то это тоже приведёт к ошибке, т.к. линковщик не сможет
%определить какую из них использовать. 

\begin{frame}[fragile]{Компилиция этап 3: линковка (компоновка)}
    \begin{itemize}
        \item На этапе компоновки важно, что каждая функция имеет уникальное имя.
        \item В \langcpp может быть две функции с одним именем, но разными
            параметрами.
        \item Имена функций искажаются (mangle) таким образом, что в их имени
            кодируются их параметры.\\
            Например, компилятор GCC превратит имя функции \verb!foo!
            \begin{lstlisting}
void foo(int, double) {}
            \end{lstlisting}
            в \verb!_Z3fooid!.
\item Аналогично функциям в линковке нуждаются глобальные переменные.
    \end{itemize}
\end{frame}

\begin{frame}[fragile]{Общая схема компиляции}
    \begin{center}
        \begin{tikzpicture}[auto,thick,font=\tiny]
            \tikzstyle{arrow} = [>=stealth',shorten >=2pt, shorten <=2pt]
            \tikzstyle{file} = [rectangle,minimum width=7ex,minimum height=15mm,draw,thick]

            \node [file,shift={(0mm,0cm)}]    (file1) {{\tt file1.cpp}};
            \node [file,shift={(0mm,-2cm)}]   (file2) {{\tt file2.cpp}};
            \node [shift={(0mm,-35mm)}]        (filed) {$\vdots$};
            \node [file,shift={(0mm,-50mm)}]   (filen) {{\tt filen.cpp}};

            \node [file,shift={(4cm,0cm)}]    (file1o) {{\tt file1.o}};
            \node [file,shift={(4cm,-2cm)}]   (file2o) {{\tt file2.o}};
            \node [     shift={(4cm,-35mm)}]  (filedo) {$\vdots$};
            \node [file,shift={(4cm,-50mm)}]  (fileno) {{\tt filen.o}};

            \node [file,shift={(8cm,-25mm)}]    (program) {{\tt program}};

            \path [arrow,->] (file1) edge node {Компиляция}  (file1o);
            \path [arrow,->] (file2) edge node {}  (file2o);
            \path [arrow,->] (filen) edge node {}  (fileno);

            \path [arrow,->] (file1o) edge node {Линковка}  (program);
            \path [arrow,->] (file2o) edge node {}  (program);
            \path [arrow,->] (fileno) edge node {}  (program);
        \end{tikzpicture}
    \end{center}
\end{frame}

\begin{frame}[fragile]{Простые типы данных С++}
    \begin{itemize}
        \item Целочисленные:
            \begin{enumerate}
                \item \code{char} (символьный тип данных)
                \item \code{short int}
                \item \code{int}
                \item \code{long int}
            \end{enumerate}
        Могут быть беззнаковыми (\code{unsigned}).
        \begin{itemize}
            \item $-2^{n-1}\ldots (2^{n-1} - 1)$ ($n$~--- число бит)
            \item $0\ldots (2^{n} - 1)$ для $\code{unsigned}$
        \end{itemize}

        \item Числа с плавающей точкой:
            \begin{enumerate}
                \item \code{float}, 4 байта, 7 значащих цифр.
                \item \code{double}, 8 байт, 15 значащих цифр.
            \end{enumerate}
        \item Логический тип данных \code{bool}.

        \item Пустой тип \code{void}.
    \end{itemize}
\end{frame}

\begin{frame}[fragile]{Замечания касательно простых типов}
    \begin{itemize}
        \item Размеры многих типов ограничены только снизу так, 
            например, тип int может быть размера 4 и 8 байт в зависимости от целевой платформы.
        \item Существуют алиасы, гарантирующие размер типа,
        они расположены в заголовочном файле \code{<cstdint>}:
            \begin{enumerate}
                \item \code{std::int32\_t}~--- знаковый, размера 32 бита (4 байта),
                \item \code{std::uint32\_t}~--- беззнаковый, размера 32 бита (4 байта),
                \item \code{std::int16\_t}~--- знаковый, размера 16 бит (2 байта),
                \item \code{std::int8\_t}~--- знаковый, размера 8 бит (1 байта),
                \item \code{std::int64\_t}~--- знаковый 64 бита (8 байт), 
                \item и т.д.
            \end{enumerate}
    \end{itemize}
\end{frame}

\begin{frame}[fragile]{Замечания касательно простых типов}
    \begin{itemize}
        \item Для получения размера конкретного типа можно использовать оператор \code{sizeof}.
        \item \code{sizeof(std::int32\_t) == 4}.
        \item Для получения максимального, минимального значения,
            можно использовать шаблон из заголовочного файла \code{<limits>}. 
        \item Например:
            \begin{enumerate}
                \item \code{std::numeric\_limits<std::uint32\_t>::max()} $= 2^{32} - 1$
                \item \code{std::numeric\_limits<std::int32\_t>::max()} $= 2^{31} - 1$
            \end{enumerate}
    \end{itemize}
\end{frame}

\begin{frame}[fragile]{Литералы}
\begin{itemize}
    \item Целочисленные:
        \begin{enumerate}
            \item \code{'a'}~--- код буквы 'a', тип \code{char},
            \item \code{42} ~--- все целые числа по умолчанию типа \code{int},
            \item \code{1234567890L}~--- суффикс '\code{L}' соответствует типу \code{long},
            \item \code{1703U}~--- суффикс '\code{U}' соответствует типу \code{unsigned int},
            \item \code{2128506UL}~--- соответствует типу \code{unsigned long}.
        \end{enumerate}

    \item Числа с плавающей точкой:
        \begin{enumerate}
            \item \code{3.14}~--- 
                все числа с точкой по умолчанию типа \code{double},
            \item \code{2.71F}~--- суффикс '\code{F}' соответствует типу 
                \code{float},
            \item \code{3.0E8}~--- соответствует $3.0\cdot 10^{8}$.
        \end{enumerate}
    \item \code{true} и \code{false}~--- значения типа \code{bool}.

    \item Строки задаются в двойных кавычках: \verb!"Text string"!.

\end{itemize}
\end{frame}

\begin{frame}[fragile]{Переменные}
\begin{itemize}
    \item При определении переменной указывается её тип. При определении можно
        сразу задать начальное значение (инициализация).
        \begin{lstlisting}
int     i = 10;
short   j = 20;
bool    b = false;

unsigned long l = 123123;

double x = 13.5, y = 3.1415;
float z;
        \end{lstlisting}
    \item Нужно всегда инициализировать переменные.
    \item Нельзя определить переменную пустого типа \code{void}.
\end{itemize}
\end{frame}

\begin{frame}[fragile]{Операции}
    \begin{itemize}
        \item Оператор присваивания: \code{=}.
        \item Арифметические:
        \begin{enumerate}
            \item бинарные: \code{+ - * / \%},
            \item унарные:  \code{++ {-}{-}}.
        \end{enumerate}
        \item Логические:
        \begin{enumerate}
            \item бинарные:  \code{\&\& ||},
            \item унарные:   \code{!}. 
        \end{enumerate}
        \item Сравнения: \code{== != > < >= <= <=>}.
        \item С-style приведение типов: \code{(type)}.
        \item Сокращённые версии бинарных операторов:
            \code{+= -= *= /= \%=}.
    \end{itemize}
\end{frame}

\begin{frame}[fragile]{Операции}
    \begin{lstlisting}
int i = 10; // initialize i with integer value
i = (20 * 3) % 7; // integer calculations

int k = i++; // post increment
int l = --i; // pre decrement

bool b = !(k == l); // logical expression
// another logical expression
b = (a == 0) || (1 / a < 1);
auto p = k <=> l; // ordereing since C++20

double d = 3.1415; // initialize d with integer value
float f = (int)d; // cast d to float
// floating point calculations
d *= i + k; // equal to d = d * (i + k)
    \end{lstlisting}
\end{frame}

\begin{frame}[fragile]{Инструкции}
    \begin{itemize}
        \item Выполнение состоит из последовательности {\em инструкций}.
        \item Инструкции выполняются одна за другой.
        \item Порядок вычислений внутри инструкций не определён.
            \begin{lstlisting}
int i = 10;
i = (i += 5) + (i * 4); // unspecified behavior
            \end{lstlisting}
        \item Блоки имеют вложенную область видимости:
            \begin{lstlisting}
int k = 10;                
{   
    int k = 5 * i;  // не видна за пределами блока
    i = (k += 5) + 5; 
}
k = k + 1;
            \end{lstlisting}

    \end{itemize}
\end{frame}

\begin{frame}[fragile]{Замечания касательно инструкций}
    \begin{itemize}
        \item В \langcpp существует несколько "особых состояний" программы,
            которые описываются в стандарте.
        \item Undefined behaviour~--- поведение программы не определено (может произойти все что угодно).
            Например, переполнение знакового целого числа.
            \begin{lstlisting}
int i = std::numeric_limits<int>::max() + 1; // UB
            \end{lstlisting}
    \end{itemize}
\end{frame}

\begin{frame}[fragile]{Замечания касательно инструкций}
    \begin{itemize}
        \item Unspecified behaviour~--- в определённых маргинальных ситуациях программа может выдавать результат, 
            зависящий от реализации компилятора.
            \begin{lstlisting}
void foo(int a, int b, int c) {
    std::cout << a << b << c; // prints a, b ,c
}

int i = 0;
foo(++i, i++, i); // unspecified
            \end{lstlisting}
    \end{itemize}
\end{frame}

\begin{frame}[fragile]{Условные операторы}
    \begin{itemize}    
        \item Оператор \code{if}:
            \begin{lstlisting}
int d = b * b - 4 * a * c;
if (d > 0) {
    roots = 2;
} else if (d == 0) {  
    roots = 1;
} else {
    roots = 0;
}
            \end{lstlisting}
        \item Тернарный условный оператор:
            \begin{lstlisting}
int roots = 0;
if (d >= 0)
    roots = (d > 0 ) ? 2 : 1;
            \end{lstlisting}
    \end{itemize}
\end{frame}

\begin{frame}[fragile]{Циклы}
    \begin{itemize}    
        \item Цикл \code{while}:
            \begin{lstlisting}
int squares = 0;
int k = 0; 
while (k < 10) {
    squares += k * k;
    k = k + 1;
} 
            \end{lstlisting}

        \item Цикл \code{for}:
            \begin{lstlisting}
for (int k = 0; k < 10; k = k + 1) {
    squares += k * k;
} 
            \end{lstlisting}
        \item Для выхода из цикла используется оператор \code{break}.
    \end{itemize}
\end{frame}

\begin{frame}[fragile]{Функции}
    \begin{itemize}
        \item В сигнатуре функции указывается тип возвращаемого значений и типы параметров.
        \item Ключевое слово \code{return} возвращает значение.
            \begin{lstlisting}
double square(double x) {
    return x * x;
}
            \end{lstlisting}
        \item Переменные, определённые внутри функций,~— {\em локальные}.
        \item Функция может возвращать \code{void}.
        \item Параметры передаются по значению (копируются).
            \begin{lstlisting}
void strange(double x, double y) {
    x = y;
}
            \end{lstlisting}
    \end{itemize}
\end{frame}

\begin{frame}[fragile]{Макросы}
    \begin{itemize}
        \item Макросами в \langcpp называют инструкции препроцессора.
        \item Препроцессор \langcpp является самостоятельным языком,
            работающим с произвольными строками.
        \item Макросы можно использовать для определения функций:
            \begin{lstlisting}
int max1(int x, int y) { 
    return x > y ? x : y;
}

#define max2(x, y)   x > y ? x : y

a = b + max2(c, d); //  b + c > d ? c : d;
            \end{lstlisting}
        \item Препроцессор ``не знает'' про синтаксис \langcpp.
    \end{itemize}
\end{frame}
    
\begin{frame}[fragile]{Макросы}
    \begin{itemize}
        \item Параметры макросов нужно оборачивать в скобки:
            \begin{lstlisting}
#define max3(x, y)   ((x) > (y) ? (x) : (y))
            \end{lstlisting}
        \item Это не избавляет от всех проблем:
            \begin{lstlisting}
int a = 1;
int b = 1;
int c = max3(++a, b); 
// c = ((++a) > (b) ? (++a) : (b))
            \end{lstlisting}
        \item Определять константы через макросы~--- плохая идея.
        \item Определять функции через макросы~--- плохая идея.
        \item Макросы можно использовать для условной компиляции:
            \begin{lstlisting}
#ifdef DEBUG
    // дополнительные проверки
#endif
            \end{lstlisting}
    \end{itemize}
\end{frame}

\begin{frame}[fragile]{Ввод-вывод}
    \begin{itemize}    
        \item Для консоли будем использовать библиотеку \code{<iostream>}.
        \begin{lstlisting}
#include <iostream>
using namespace std;
            \end{lstlisting}
        \item Ввод:
            \begin{lstlisting}
int a = 0;
int b = 0;
cin >> a >> b;
            \end{lstlisting}
        \item Вывод:
            \begin{lstlisting}
cout << "a + b = " << (a + b) << endl;
            \end{lstlisting}
        \item Аналогичным образов осуществляется ввод-вывод из файла \code{<fstream>}
    \end{itemize}
\end{frame}

\begin{frame}[fragile]{Простая программа}
    \begin{lstlisting}
#include <iostream>
using namespace std;

int main () 
{
    int a = 0;
    int b = 0;

    cout << "Enter a and b: ";
    cin >> a >> b;

    cout << "a + b = " << (a + b) << endl;

    return 0;
}
    \end{lstlisting}
\end{frame}

\begin{frame}[fragile]{Контракты при разработке программ}
    \begin{itemize}
        \item \code{assert(<expression>)} или \code{static\_assert(<expression>, <error-text>)} 
            позволяют выполнять проверки во время исполенения и компиляции соотвественно.
        \item Такие проверки удобны во время разработки больших программных продуктов и рефакторинга.
        \item Важно отметить, что код для \code{assert(<expression>)} 
            по умолчанию исполняется только в debug сборке.
    \end{itemize}
\end{frame}

\begin{frame}[fragile]{Архитектура фон Неймана}
    Современных компьютеры построены по принципам 
    архитектуры фон Неймана:
    \begin{enumerate}
        \item {\bf Принцип однородности памяти.}\\
            Команды и данные хранятся в одной и той же памяти и внешне в памяти неразличимы.
        \item {\bf Принцип адресности.}\\
            Память состоит из пронумерованных ячеек.
        \item {\bf Принцип программного управления.}\\
            Все вычисления представляются в виде последовательности команд.
        \item {\bf Принцип двоичного кодирования.}\\
            Вся информация (данные и команды) кодируются двоичными числами.
    \end{enumerate}
\end{frame}

\begin{frame}[fragile]{Сегментация памяти}
    \begin{itemize}
        \item Оперативная память, используемая в программе на C++,
            разделена на области двух типов:
            \begin{enumerate}
                \item сегменты данных,
                \item сегменты кода (текстовые сегменты).
            \end{enumerate}
        \item В сегментах кода содержится код программы.
        \item В сегментах данных располагаются данные программы 
            (значения переменных, массивы и пр.).
        \item При запуске программы выделяются два сегмента данных:
            \begin{enumerate}
                \item сегмент глобальных данных,
                \item стек.
            \end{enumerate}
        \item В процессе работы программы могут выделяться и освобождаться
            дополнительные сегменты памяти.
        \item Обращения к адресу вне выделенных сегментов~--- ошибка времени
            выполнения (access violation, segmentation fault).
    \end{itemize}
\end{frame}

\begin{frame}[fragile]{Как выполняется программа?}
    \begin{itemize}
        \item Каждой функции в скомпилированном коде соответствует отдельная
            секция.
        \item Адрес начала такой секции — это адрес функции.
        \item Телу функции соответствует последовательность команд процессора.
        \item Работа с данными происходит на уровне байт,
            информация о типах отсутствует.
        \item В процессе выполнения адрес следующей инструкции хранится 
            в специальном регистре процессора IP (Instruction Pointer).
        \item Команды выполняются последовательно, пока не встретится
            специальная команда (например, условный переход или вызов 
            функции), которая изменит IP.
    \end{itemize}
\end{frame}

\begin{frame}[fragile]{Ещё раз о линковке}
    \begin{itemize}
        \item На этапе компиляции объектных файлов
            в места вызова функций подставляются имена функций.
        \item На этапе линковки в места вызова вместо имён функций
            подставляются их адреса.
        \item Ошибки линковки:
            \begin{enumerate}
                \item {\tt undefined reference}\\
                    Функция имеет объявление, но не имеет тела.

                \item {\tt multiple definition}\\
                    Функция имеет два или более определений.
            \end{enumerate}
        \item Наиболее распространённый способ получить 
            {\tt multiple definition}~--- определить функцию
            в заголовочном файле, который включён в несколько
            {\tt .cpp} файлов.
    \end{itemize}
\end{frame}

\begin{frame}[fragile]{Стек вызовов}
    \begin{itemize}
        \item Стек вызовов~— это сегмент данных, используемый для хранения локальных
            переменных и временных значений.
        \item Не стоит путать стек с одноимённой структурой данных, у~стека в \langcpp
            можно обратиться к произвольной ячейке.
        \item Стек выделяется при запуске программы.
        \item Стек обычно небольшой по размеру (4Мб).
        \item Функции хранят свои локальные переменные на стеке. 
        \item При выходе из функции соответствующая область стека
            объявляется свободной.
        \item Промежуточные значения, возникающие при вычислении
            сложных выражений, также хранятся на стеке.
    \end{itemize}
\end{frame}

\begin{frame}[fragile]{Устройство стека}
    \begin{minipage}{.5\textwidth}
        \begin{tikzpicture}[
            start chain=1 going below,node distance=-0.15mm
            ]
            \tikzstyle{cell} = [rectangle,minimum width=2cm,draw,thick,on chain=1];
            \tikzstyle{function} = [cell, minimum height=1cm]; 
            \tikzstyle{arrow} = [>=stealth',shorten >=2pt, shorten <=2pt]

            \node [function, fill=green!30]  (main) {main()};
            \node<1,5> [cell, minimum height=45mm, dashed] (empty) {};

            \node<2,4> [function, fill=blue!30]   (foo)  {foo()};
            \node<2,4> [cell, minimum height=35mm, dashed] (empty) {};

            \node<3> [function, fill=blue!30] (foo)  {foo()};
            \node<3> [function, fill=orange!30] (bar)  {bar()};
            \node<3> [cell, minimum height=25mm, dashed] (empty) {};

            \node<6> [function, fill=orange!30] (bar)  {bar()};
            \node<6> [cell, minimum height=35mm, dashed] (empty) {};

            \node [right of=main, xshift=3cm, yshift=8mm] (bottom) {дно стека};
            \path [arrow,->] (bottom) edge node {} (main.north east);

            \node [right of=empty, xshift=3cm, yshift=-8mm] (top) {вершина стека};
            \path [arrow,->] (top) edge node {} (empty.north east);
        \end{tikzpicture}
    \end{minipage}\hfill
    \begin{minipage}{.4\textwidth}
        \begin{lstlisting}
void bar( ) {
    int c;
}

void foo( ) {
    int b = 3;
    bar();
}

int main( ) {
    int a = 3;    
    foo();
    bar();

    return 0;
}
        \end{lstlisting}
    \end{minipage}
\end{frame}
    
\begin{frame}[fragile]{Вызов функции}
    \begin{minipage}{.3\textwidth}
    \small
        \begin{tikzpicture}[
            start chain=1 going below,node distance=-0.15mm
            ]
            \tikzstyle{cell} = [rectangle,minimum width=2cm, minimum height=5mm,,draw,thick,on chain=1];
            \tikzstyle{function} = [cell, minimum height=1cm]; 
            \tikzstyle{arrow} = [>=stealth',shorten >=2pt, shorten <=2pt]
            \tikzstyle{mbrace} = [decorate,decoration={brace,amplitude=5pt}];
        
            \node<1-8> [cell, fill=green!30] (x)   {\tt x = 1};
            \node<9>   [cell, fill=red!30]   (x)   {\tt x = 2};
            \node [cell, fill=green!30] (y)   {\tt y = 2};
        
        
            \node<2-8> [cell, fill=green!30] (p3)   {\tt false};
            \node<2-8> [cell, fill=green!30] (p2)   {\tt 2};
            \node<2-8> [cell, fill=green!30] (p1)   {\tt 1};
        
        %    \draw[mbrace] 
        %        (p3.north east)--(p1.south east) node[midway,anchor=west,xshift=6pt]
        %        {аргументы};
        
            \node<3-6> [cell, fill=green!30] (retv)  {ret val};
            \node<7-8> [cell, fill=red!30]   (retv)  {2};
            \node<3-8> [cell, fill=green!30] (reta)  {ret addr};
            \node<3-8> [cell, fill=green!30] (retr)  {registers};
        
        
            %\draw<1-3>[mbrace] 
            %    (x.north east)--(empty.north east) node[midway,anchor=west,xshift=6pt] {main};
        
        %    \node[right of=retv,xshift=21mm] {результат};
        %    \node[right of=reta,xshift=15mm] {IP};
        %    \node[right of=retr,xshift=21mm] {состояние};


            \node<4-5> [cell, fill=blue!30] (d)   {\tt d};
            \node<4-5> [cell, fill=blue!30] (h)   {\tt h};
            \node<5>   [cell, fill=red!30] (tmp)  {\tt a * b = 2};
            \node<6-7> [cell, fill=blue!30] (d)   {\tt d = 5.42};
            \node<6-7> [cell, fill=blue!30] (h)   {\tt h = 2};
        
        
            \node<1,9> [cell, minimum height=55mm, dashed] (empty) {};
            \node<2> [cell, minimum height=40mm, dashed] (empty) {};
            \node<3,8> [cell, minimum height=25mm, dashed] (empty) {};
            \node<4> [cell, minimum height=15mm, dashed] (empty) {};
            \node<5> [cell, minimum height=10mm, dashed] (empty) {};
            \node<6-7> [cell, minimum height=15mm, dashed] (empty) {};

            \node [right of=y,xshift=22mm] (fp) {\parbox{15mm}{frame\\ pointer}};
            
            \path<1-3,8-9> [arrow,->] (fp) edge node {} (x.north east);

            \path<4-7> [arrow,->] (fp) edge node {} (d.north east);

            %\draw<4> [mbrace] 
            %    (d.north east)--(h.south east) node[midway,anchor=west,xshift=8pt] {foo};
                                            
            \node [right of=empty,xshift=22mm] (sp) {\parbox{15mm}{stack\\ pointer}};
            \path [arrow,->] (sp) edge node {} (empty.north east);
        \end{tikzpicture}
    \end{minipage}\hfill
    \begin{minipage}{.6\textwidth}
        \begin{lstlisting}
int foo(int a, int b, bool c) 
{
    double d = a * b * 2.71;
    int    h = c ? d : d / 2;
    return h;
}

int main( ) 
{
    int x = 1;
    int y = 2;
    x = foo (x, y, false);
    cout << x;
    return 0;
}
        \end{lstlisting}
    \end{minipage}
\end{frame}

\begin{frame}[fragile]{Вызов функции}
    \begin{itemize}
        \item При вызове функции на стек складываются:
            \begin{enumerate}
                \item aргументы функции,
                \item адрес возврата,
                \item значение frame pointer и регистров процессора.
            \end{enumerate}
        \item Кроме этого на стеке резервируется место под
            возвращаемое значение.
        \item Параметры передаются в обратном порядке, что позволяет
            реализовать функции с переменным числом аргументов.
        \item Адресация локальных переменных функции и аргументов
            функции происходит относительно frame pointer.
        \item Конкретный процесс вызова зависит от используемых 
            соглашений (cdecl,  stdcall, fastcall, thiscall).
    \end{itemize}
\end{frame}

\end{document}

